\documentclass[12pt, a4paper, openany]{report}

% ── Encodage & langue ────────────────────────────────────────────────────────
\usepackage[utf8]{inputenc}
\usepackage[T1]{fontenc}
\usepackage[french]{babel}

% ── Mise en page ─────────────────────────────────────────────────────────────
\usepackage[top=2.5cm, bottom=2.5cm, left=2.8cm, right=2.8cm]{geometry}
\usepackage{setspace}
\onehalfspacing

% ── Polices & couleurs ───────────────────────────────────────────────────────
\usepackage{lmodern}
\usepackage[dvipsnames]{xcolor}
\definecolor{primary}{RGB}{124, 58, 237}
\definecolor{secondary}{RGB}{91, 33, 182}
\definecolor{lightgray}{RGB}{245, 245, 245}
\definecolor{medgray}{RGB}{210, 210, 210}
\definecolor{darkgray}{RGB}{70, 70, 70}
\definecolor{success}{RGB}{34, 197, 94}
\definecolor{warning}{RGB}{251, 146, 60}
\definecolor{danger}{RGB}{239, 68, 68}

% ── Graphiques & figures ─────────────────────────────────────────────────────
\usepackage{graphicx}
\usepackage{float}
\usepackage{caption}
\usepackage{subcaption}
\usepackage{tikz}
\usetikzlibrary{shapes.geometric, arrows.meta, positioning, shadows, calc}

% ── Tableaux ─────────────────────────────────────────────────────────────────
\usepackage{booktabs}
\usepackage{tabularx}
\usepackage{array}
\usepackage{multirow}
\usepackage{colortbl}

% ── Listes & énumérations ────────────────────────────────────────────────────
\usepackage{enumitem}
\usepackage{pifont}
\newcommand{\cmark}{\textcolor{success}{\ding{51}}}
\newcommand{\xmark}{\textcolor{danger}{\ding{55}}}
\newcommand{\pmark}{\textcolor{warning}{\ding{119}}}

% ── Code source ──────────────────────────────────────────────────────────────
\usepackage{listings}
\lstset{
  basicstyle=\ttfamily\footnotesize,
  backgroundcolor=\color{lightgray},
  frame=single,
  rulecolor=\color{medgray},
  breaklines=true,
  showstringspaces=false,
  keywordstyle=\color{primary}\bfseries,
  commentstyle=\color{darkgray}\itshape,
  stringstyle=\color{secondary},
  numbers=left,
  numberstyle=\tiny\color{darkgray},
  tabsize=4,
  xleftmargin=1.5em,
  framexleftmargin=1.5em
}

% ── Liens & PDF ──────────────────────────────────────────────────────────────
\usepackage{hyperref}
\hypersetup{
  colorlinks=true,
  linkcolor=primary,
  urlcolor=secondary,
  citecolor=primary,
  pdftitle={Rapport de mi-parcours -- TravelShare},
  pdfauthor={Mahmoud}
}

% ── Titres de chapitres et sections ──────────────────────────────────────────
\usepackage{titlesec}

% Chapitre : numéro dans une boîte colorée + titre en violet + filet
\titleformat{\chapter}[hang]
  {\Huge\bfseries}
  {\colorbox{primary}{\raisebox{0pt}[16pt][4pt]{\color{white}\quad\thechapter\quad}}\enspace}
  {0pt}
  {\Huge\bfseries\color{primary}}
  [\vspace{6pt}{\color{primary!30}\titlerule[1.5pt]}]
\titlespacing*{\chapter}{0pt}{0pt}{20pt}

% Section : numéro + filet fin
\titleformat{\section}
  {\Large\bfseries\color{primary}}
  {\thesection\enspace\textbar\enspace}
  {0pt}
  {}
  [\vspace{-2pt}{\color{primary!20}\titlerule[0.6pt]}]
\titlespacing*{\section}{0pt}{14pt}{6pt}

% Sous-section
\titleformat{\subsection}
  {\large\bfseries\color{darkgray}}
  {\thesubsection\enspace}
  {0pt}
  {}
\titlespacing*{\subsection}{0pt}{10pt}{4pt}

% ── En-têtes et pieds de page ────────────────────────────────────────────────
\usepackage{fancyhdr}
\pagestyle{fancy}
\fancyhf{}
\renewcommand{\headrulewidth}{1.2pt}
\renewcommand{\headrule}{{\color{primary}\hrule width\headwidth height\headrulewidth}}
\fancyhead[L]{\small\textcolor{primary}{\textbf{TravelShare}}\;\textcolor{medgray}{|}\;\textcolor{darkgray}{\textit{Rapport de mi-parcours}}}
\fancyhead[R]{\small\textcolor{darkgray}{\nouppercase{\leftmark}}}
\fancyfoot[C]{\small\textcolor{darkgray}{\thepage}}
\fancypagestyle{plain}{%
  \fancyhf{}
  \renewcommand{\headrulewidth}{0pt}
  \fancyfoot[C]{\small\textcolor{darkgray}{\thepage}}
}

% ── Boîtes colorées ──────────────────────────────────────────────────────────
\usepackage{tcolorbox}
\tcbuselibrary{skins, breakable}

\newtcolorbox{infobox}[1][]{
  colback=primary!4!white, colframe=primary,
  fonttitle=\bfseries, title=#1,
  boxrule=0.6pt, arc=4pt, left=6pt, right=6pt, top=4pt, bottom=4pt
}

\newtcolorbox{warningbox}[1][]{
  colback=orange!7!white, colframe=warning,
  fonttitle=\bfseries, title=#1,
  boxrule=0.6pt, arc=4pt, left=6pt, right=6pt, top=4pt, bottom=4pt
}

\newtcolorbox{successbox}[1][]{
  colback=green!4!white, colframe=success,
  fonttitle=\bfseries, title=#1,
  boxrule=0.6pt, arc=4pt, left=6pt, right=6pt, top=4pt, bottom=4pt
}

% ── Commandes personnalisées ──────────────────────────────────────────────────
\newcommand{\placeholder}[2]{%
  \begin{figure}[H]
    \centering
    \begin{tikzpicture}
      \node[draw=medgray, dashed, fill=lightgray, minimum width=#1, minimum height=3cm,
            font=\small\itshape\color{darkgray}, align=center, rounded corners=4pt] {#2};
    \end{tikzpicture}
    \caption{#2}
  \end{figure}
}
\newcommand{\todo}[1]{\textcolor{danger}{\textbf{[TODO: #1]}}}

% ─────────────────────────────────────────────────────────────────────────────
\begin{document}

% ── PAGE DE GARDE ─────────────────────────────────────────────────────────────
\begin{titlepage}
\newgeometry{top=0cm, bottom=0cm, left=0cm, right=0cm}
\thispagestyle{empty}

% Bande supérieure
\noindent
\begin{tikzpicture}[remember picture, overlay]
  \fill[primary] (current page.north west)
    rectangle ([yshift=-4.2cm]current page.north east);
  \node[anchor=west, text=white, font=\large\bfseries]
    at ([xshift=1.8cm, yshift=-1.1cm]current page.north west)
    {Université de Montpellier};
  \node[anchor=west, text=white!75!primary, font=\normalsize]
    at ([xshift=1.8cm, yshift=-1.9cm]current page.north west)
    {Faculté des Sciences --- Département Informatique};
  \node[anchor=west, text=white!75!primary, font=\normalsize]
    at ([xshift=1.8cm, yshift=-2.7cm]current page.north west)
    {Année universitaire 2025--2026};
\end{tikzpicture}

\vspace{5cm}

% Zone centrale
\begin{center}

  % Logo TravelShare
  \begin{tikzpicture}
    \node[fill=primary!10, draw=primary, line width=1.5pt,
          minimum width=9cm, minimum height=1.6cm,
          rounded corners=8pt,
          font=\LARGE\bfseries\color{primary}]
      {\texttt{travel}\textbf{$\bullet$}\texttt{share}};
  \end{tikzpicture}

  \vspace{1.4cm}
  {\color{medgray}\rule{13cm}{0.6pt}}
  \vspace{1cm}

  {\fontsize{30}{36}\selectfont\bfseries\color{primary} Rapport de mi-parcours}

  \vspace{0.6cm}
  {\Large\color{darkgray} Projet Android \textemdash{} \textit{Traveling}}

  \vspace{0.4cm}
  {\large\color{secondary}
    Partie~1~: TravelShare \textemdash{} Frontend Android \& Backend REST}

  \vspace{1.2cm}
  {\color{medgray}\rule{13cm}{0.6pt}}
  \vspace{1cm}

  \renewcommand{\arraystretch}{1.6}
  \begin{tabular}{r@{\hspace{1cm}}l}
    \textcolor{darkgray}{\textbf{Auteur}}         & Mahmoud \\
    \textcolor{darkgray}{\textbf{Module}}         & Développement Android \\
    \textcolor{darkgray}{\textbf{Dépôt Android}}  & \texttt{travelshare-android} \\
    \textcolor{darkgray}{\textbf{Dépôt Backend}}  & \texttt{travelshare-backend} \\
  \end{tabular}

\end{center}

\vfill

% Bande inférieure
\begin{tikzpicture}[remember picture, overlay]
  \fill[primary] (current page.south west)
    rectangle ([yshift=2cm]current page.south east);
  \node[text=white, font=\normalsize\bfseries]
    at ([yshift=1.3cm]current page.south)
    {27 avril 2026};
  \node[text=white!75!primary, font=\small]
    at ([yshift=0.6cm]current page.south)
    {Application Android de partage et découverte de photos de voyages};
\end{tikzpicture}

\restoregeometry
\end{titlepage}

% ── TABLE DES MATIÈRES ────────────────────────────────────────────────────────
\tableofcontents
\newpage

% ─────────────────────────────────────────────────────────────────────────────
\chapter{Introduction et contexte du projet}
% ─────────────────────────────────────────────────────────────────────────────

\section{Présentation générale}

Le projet \textbf{Traveling} est une application mobile Android destinée aux voyageurs. Il est
divisé en deux parties principales développées en binôme~:

\begin{itemize}
  \item \textbf{TravelShare (Partie~1)} : application de partage et de découverte de photos
        de voyages -- dont la réalisation m'est confiée.
  \item \textbf{TravelPath (Partie~2)} : générateur intelligent de parcours de visite,
        développé par le second membre du binôme.
  \item \textbf{Passerelle (Partie~3)} : intégration cohérente des deux modules en une
        seule application unifiée (\textit{Traveling}).
\end{itemize}

Ce rapport couvre l'état d'avancement de la \textbf{partie TravelShare} à mi-parcours du
développement.

\section{Objectifs de TravelShare}

TravelShare est une application permettant à des voyageurs de partager, découvrir et interagir
autour de photos de voyage. Elle propose deux modes de fonctionnement~:

\begin{infobox}[Mode anonyme]
Sans compte utilisateur, le visiteur peut~: parcourir les photos publiques, effectuer des
recherches (textuelles ou vocales), consulter les détails d'une photo (lieu, date,
commentaires, itinéraire), aimer ou signaler un contenu.
\end{infobox}

\begin{infobox}[Mode connecté]
Après inscription et connexion, l'utilisateur peut en plus~: publier des photos, les partager
avec des groupes ou les rendre publiques, enrichir les photos (texte, tags, annotations IA),
recevoir des notifications et accéder à un profil personnel.
\end{infobox}

% ─────────────────────────────────────────────────────────────────────────────
\chapter{Architecture technique -- Frontend Android}
% ─────────────────────────────────────────────────────────────────────────────

\section{Choix technologiques}

L'application est développée en \textbf{Java} natif pour Android, en respectant les bonnes
pratiques recommandées par Google.

\subsection{Patron de conception MVVM}

L'architecture retenue est \textbf{Model-View-ViewModel (MVVM)}, soutenue par les bibliothèques
Jetpack de Google~:

\begin{itemize}
  \item \textbf{ViewModel} : séparation logique métier / couche d'affichage, survie aux
        rotations d'écran.
  \item \textbf{LiveData} : observabilité réactive des données, mise à jour automatique de
        l'interface.
  \item \textbf{Fragment} : navigation modulaire (HomeFragment, SearchFragment, ProfileFragment).
\end{itemize}

\subsection{Bibliothèques utilisées}

\begin{table}[H]
\centering
\caption{Principales dépendances du projet Android}
\begin{tabularx}{\textwidth}{lXl}
\toprule
\rowcolor{primary!15}
\textbf{Bibliothèque} & \textbf{Rôle} & \textbf{Version} \\
\midrule
Retrofit2 + OkHttp      & Client HTTP REST, communication avec le backend & 2.x \\
Gson Converter          & Désérialisation JSON $\leftrightarrow$ Java POJO & -- \\
Glide                   & Chargement et mise en cache d'images & 4.16.0 \\
Google Maps SDK         & Affichage carte, marqueurs, sélection de lieu & 19.0.0 \\
Firebase Messaging      & Notifications push (FCM) & BOM \\
PhotoView               & Zoom sur image en plein écran & 2.3.0 \\
Cloudinary              & Hébergement et upload d'images dans le cloud & -- \\
Material Design         & Composants UI (BottomSheet, ChipGroup, Badge) & -- \\
FusedLocationProvider   & Géolocalisation de l'appareil & -- \\
\bottomrule
\end{tabularx}
\end{table}

\subsection{Communication réseau}

Le client HTTP est encapsulé dans la classe \texttt{RetrofitClient}, configurée avec~:
\begin{itemize}
  \item Une URL de base dynamique (modifiable depuis l'écran \texttt{IpConfigActivity}
        pour faciliter les tests en réseau local).
  \item Un intercepteur OkHttp qui injecte automatiquement le token JWT (\texttt{Bearer})
        dans chaque requête sortante lorsque l'utilisateur est connecté.
  \item Des délais d'expiration de 60~secondes (connexion, lecture, écriture).
  \item Un intercepteur de journalisation pour le débogage.
\end{itemize}

\subsection{Structure des packages}

\begin{lstlisting}[language=Java, caption=Organisation des packages Java]
com.example.travelshare/
+-- api/                    # Ancien client API (remplace par network/)
+-- data/                   # Donnees mock (tests)
+-- model/                  # POJOs (Photo, Group, Comment, Notification)
|   +-- response/           # Objets de reponse API (AuthResponse, etc.)
+-- network/                # RetrofitClient, ApiService, TokenManager, FCM
+-- ui/
|   +-- auth/               # AuthDialog (connexion / inscription)
|   +-- detail/             # PhotoDetailActivity, PhotoDetailViewModel
|   +-- home/               # HomeFragment, adaptateurs, viewmodels, dialogs
|   +-- profile/            # ProfileFragment
|   +-- search/             # SearchFragment, SearchViewModel
+-- IpConfigActivity.java   # Configuration IP du serveur
+-- MainActivity.java       # Point d'entree, navigation bas de page
+-- MainViewModel.java      # Etat global (connexion, badge notifs)
+-- TravelShareApplication.java
\end{lstlisting}

% ─────────────────────────────────────────────────────────────────────────────
\chapter{Architecture et réalisation du Backend}
% ─────────────────────────────────────────────────────────────────────────────

\section{Vue d'ensemble}

Le backend de TravelShare est une \textbf{API REST} développée en \textbf{Node.js} avec le
framework \textbf{Express.js}. Il suit une architecture \textbf{MVC} (Modèle-Contrôleur)
et expose l'ensemble des données via des routes JSON consommées par l'application Android.

\begin{infobox}[Rôle du serveur]
Le serveur centralise~: l'authentification JWT, la persistance des données (MongoDB Atlas),
l'upload des images (Cloudinary), l'envoi des notifications push (Firebase Cloud Messaging)
et la gestion des groupes et des notifications in-app.
\end{infobox}

\section{Choix technologiques}

\begin{table}[H]
\centering
\caption{Dépendances du backend Node.js}
\renewcommand{\arraystretch}{1.3}
\begin{tabularx}{\textwidth}{llX}
\toprule
\rowcolor{primary!15}
\textbf{Technologie} & \textbf{Version} & \textbf{Rôle} \\
\midrule
Node.js / Express.js            & 4.18.x & Serveur HTTP, routage, middleware \\
MongoDB + Mongoose              & 7.3.x  & Base de données NoSQL, ODM \\
bcryptjs                        & 2.4.x  & Hachage des mots de passe (sel de 10 tours) \\
jsonwebtoken (JWT)              & 9.0.x  & Génération et vérification des tokens \\
Cloudinary SDK                  & 1.41.x & Upload et hébergement des images \\
multer + multer-storage-cloudinary & 2.x / 4.x & Traitement multipart, stockage Cloudinary \\
cors                            & 2.8.x  & Contrôle des requêtes cross-origin \\
dotenv                          & 16.x   & Variables d'environnement \\
nodemon                         & 3.x    & Rechargement automatique (développement) \\
\bottomrule
\end{tabularx}
\end{table}

\section{Structure du projet}

\begin{lstlisting}[language=bash, caption=Organisation du projet backend]
travelshare-backend/
+-- config/
|   +-- db.js             # Connexion MongoDB (Mongoose)
|   +-- cloudinary.js     # Configuration Cloudinary + Multer
|   +-- firebase.js       # Client FCM HTTP v1 (sans SDK)
+-- controllers/
|   +-- authController.js
|   +-- photoController.js
|   +-- groupController.js
|   +-- notificationController.js
+-- middleware/
|   +-- authMiddleware.js  # protect + optionalAuth
+-- models/
|   +-- User.js
|   +-- Photo.js
|   +-- Group.js
|   +-- Notification.js
+-- routes/
|   +-- auth.js
|   +-- photos.js
|   +-- groups.js
|   +-- notifications.js
+-- server.js              # Point d'entree, montage des routes
\end{lstlisting}

\section{Modèles de données (Mongoose)}

\subsection{Modèle \texttt{User}}

Le schéma \texttt{User} stocke les informations de compte~:

\begin{lstlisting}[language=JavaScript, caption=Schéma User.js (extrait)]
{
  fullName  : String (required),
  username  : String (unique, lowercase),
  email     : String (unique, regex validation),
  password  : String (minlength 6, select: false),
  avatar    : String,
  fcmToken  : String,   // token Firebase pour les push
  createdAt : Date
}
\end{lstlisting}

Le mot de passe est haché automatiquement via un hook \texttt{pre('save')} avec \textbf{bcryptjs}
(sel aléatoire, 10 tours). La méthode d'instance \texttt{matchPassword()} permet la vérification
à la connexion. Le champ \texttt{password} est systématiquement exclu des réponses API
(\texttt{select: false}).

\subsection{Modèle \texttt{Photo}}

\begin{lstlisting}[language=JavaScript, caption=Schéma Photo.js (extrait)]
{
  author       : ObjectId -> User (required),
  imageUrl     : String,
  description  : String,
  location     : String,       // ville
  country      : String,
  locationType : enum [Nature, Urbain, Plage, Montagne, Musee,
                        Restaurant, Magasin, Rue, Autre],
  coordinates  : GeoJSON Point, // index 2dsphere
  tags         : [String],
  howToGetThere: String,
  likes        : [ObjectId -> User],
  isPublic     : Boolean (default true),
  groups       : [ObjectId -> Group],
  reports      : [{ user, reason, date }],
  comments     : [{ user, text, createdAt }],
  createdAt    : Date
}
\end{lstlisting}

Deux index MongoDB sont définis~: un \textbf{index géospatial} \texttt{2dsphere} sur
\texttt{coordinates} (requêtes de proximité \texttt{\$centerSphere}) et un \textbf{index texte}
sur \texttt{location}, \texttt{description} et \texttt{tags} (recherche full-text
\texttt{\$text}).

\subsection{Modèle \texttt{Group}}

\begin{lstlisting}[language=JavaScript, caption=Schéma Group.js]
{
  name       : String (required),
  description: String,
  creator    : ObjectId -> User,
  members    : [ObjectId -> User],
  coverImage : String (URL),
  photos     : [ObjectId -> Photo],
  createdAt  : Date
}
\end{lstlisting}

\subsection{Modèle \texttt{Notification}}

\begin{lstlisting}[language=JavaScript, caption=Schéma Notification.js]
{
  recipient : ObjectId -> User,
  sender    : ObjectId -> User,
  type      : enum [new_photo, new_like, group_photo,
                    new_comment, group_invite],
  title     : String,
  message   : String,
  photo     : ObjectId -> Photo,
  group     : ObjectId -> Group,
  isRead    : Boolean (default false),
  relatedId : String,
  createdAt : Date
}
\end{lstlisting}

Un hook \textbf{post-save} déclenche automatiquement l'envoi d'une notification push Firebase
dès qu'une notification est persistée en base, si le destinataire possède un token FCM valide.

\section{API REST -- Référentiel des endpoints}

\begin{table}[H]
\centering
\caption{Endpoints de l'API TravelShare}
\renewcommand{\arraystretch}{1.3}
\begin{tabularx}{\textwidth}{llXc}
\toprule
\rowcolor{primary!20}
\textbf{Méthode} & \textbf{Route} & \textbf{Description} & \textbf{Auth} \\
\midrule
\multicolumn{4}{l}{\small\textit{\textbf{Authentification}}} \\
POST   & \texttt{/api/auth/register}         & Créer un compte                       & -- \\
POST   & \texttt{/api/auth/login}             & Se connecter, obtenir un JWT          & -- \\
GET    & \texttt{/api/auth/me}                & Profil de l'utilisateur connecté      & JWT \\
PUT    & \texttt{/api/auth/fcm-token}         & Enregistrer le token FCM              & JWT \\
\midrule
\multicolumn{4}{l}{\small\textit{\textbf{Photos}}} \\
GET    & \texttt{/api/photos}                 & Fil public (filtres, pagination)      & Opt. \\
GET    & \texttt{/api/photos/random}          & Flux aléatoire (10 photos)            & -- \\
GET    & \texttt{/api/photos/me}              & Mes photos publiées                   & JWT \\
POST   & \texttt{/api/photos/upload}          & Upload image vers Cloudinary          & JWT \\
GET    & \texttt{/api/photos/:id}             & Détail d'une photo                    & Opt. \\
POST   & \texttt{/api/photos}                 & Publier une nouvelle photo            & JWT \\
DELETE & \texttt{/api/photos/:id}             & Supprimer sa photo                    & JWT \\
POST   & \texttt{/api/photos/:id/toggle-like} & Liker / unliker                       & JWT \\
POST   & \texttt{/api/photos/:id/report}      & Signaler un contenu                   & JWT \\
GET    & \texttt{/api/photos/:id/comments}    & Lister les commentaires               & Opt. \\
POST   & \texttt{/api/photos/:id/comments}    & Ajouter un commentaire                & JWT \\
DELETE & \texttt{/api/photos/:id/comments/:c} & Supprimer son commentaire             & JWT \\
\midrule
\multicolumn{4}{l}{\small\textit{\textbf{Groupes}}} \\
GET    & \texttt{/api/groups}                    & Mes groupes                        & JWT \\
POST   & \texttt{/api/groups}                    & Créer un groupe                    & JWT \\
GET    & \texttt{/api/groups/discover}           & Découvrir tous les groupes         & JWT \\
PUT    & \texttt{/api/groups/:id/join}           & Rejoindre un groupe                & JWT \\
PUT    & \texttt{/api/groups/:id/leave}          & Quitter un groupe                  & JWT \\
GET    & \texttt{/api/groups/:id/photos}         & Photos partagées dans un groupe    & JWT \\
POST   & \texttt{/api/groups/:id/photos}         & Ajouter une photo au groupe        & JWT \\
DELETE & \texttt{/api/groups/:id/photos/:pid}    & Retirer une photo du groupe        & JWT \\
\midrule
\multicolumn{4}{l}{\small\textit{\textbf{Notifications}}} \\
GET    & \texttt{/api/notifications}              & Historique des notifications      & JWT \\
GET    & \texttt{/api/notifications/unread-count} & Nombre de non-lues                & JWT \\
PUT    & \texttt{/api/notifications/:id/read}     & Marquer une notification comme lue & JWT \\
PUT    & \texttt{/api/notifications/read-all}     & Tout marquer comme lu             & JWT \\
\bottomrule
\multicolumn{4}{l}{\footnotesize JWT~= token Bearer obligatoire \quad Opt.~= token optionnel (enrichit la réponse si présent)}
\end{tabularx}
\end{table}

\section{Middleware d'authentification}

Deux middlewares sont définis dans \texttt{authMiddleware.js}~:

\begin{itemize}
  \item \textbf{\texttt{protect}} : vérifie et décode le token JWT depuis l'en-tête
        \texttt{Authorization: Bearer <token>}. En cas d'échec (token absent, expiré ou
        invalide), renvoie \texttt{401 Unauthorized}. Si valide, attache l'objet utilisateur
        complet (sans \texttt{password}) à \texttt{req.user} et appelle \texttt{next()}.
  \item \textbf{\texttt{optionalAuth}} : même logique, mais en cas d'absence ou d'invalidité
        du token positionne \texttt{req.user = null} et continue. Utilisé pour les routes
        accessibles en mode anonyme (liste des photos, détail d'une photo, commentaires).
\end{itemize}

\section{Intégrations externes}

\subsection{Cloudinary -- Hébergement des images}

La configuration Cloudinary est centralisée dans \texttt{config/cloudinary.js}.
L'upload est géré via \textbf{Multer} couplé à \texttt{multer-storage-cloudinary}~: les
fichiers multipart envoyés par le client Android sont directement streamés vers le bucket
\texttt{travelshare} sur Cloudinary (formats acceptés~: \texttt{jpg}, \texttt{jpeg},
\texttt{png}, \texttt{webp}). L'URL publique retournée est stockée dans le champ
\texttt{imageUrl} du document \texttt{Photo} en base.

\subsection{Firebase Cloud Messaging -- Notifications push}

L'envoi de notifications push n'utilise pas le SDK Firebase Admin (non adapté à cet
environnement), mais appelle directement l'\textbf{API HTTP FCM v1} via le module natif
\texttt{https}. Le mécanisme comprend trois étapes~:

\begin{enumerate}
  \item Génération d'un \textbf{JWT service-account} signé avec la clé privée du compte de
        service Firebase (variable \texttt{FIREBASE\_SERVICE\_ACCOUNT} dans \texttt{.env}).
  \item Échange de ce JWT contre un \textbf{OAuth2 access token} via
        \texttt{oauth2.googleapis.com/token} (mis en cache en mémoire, renouvellement
        automatique 60~secondes avant expiration).
  \item Envoi \texttt{POST} à l'endpoint
        \texttt{fcm.googleapis.com/v1/projects/\{project\_id\}/messages:send} avec le
        token FCM du destinataire, le titre, le corps et les données complémentaires.
\end{enumerate}

Ce flux est déclenché automatiquement par le hook \texttt{post-save} du modèle
\texttt{Notification}~: tout événement persisté en base (nouveau like, commentaire, etc.)
génère immédiatement une notification push vers l'appareil du destinataire.

\section{Sécurité et bonnes pratiques}

\begin{itemize}
  \item \textbf{Mots de passe}~: hachés avec bcryptjs, jamais renvoyés dans les réponses
        (\texttt{select: false}).
  \item \textbf{JWT}~: secret externalisé en variable d'environnement
        (\texttt{JWT\_SECRET}), durée de vie paramétrable (\texttt{JWT\_EXPIRES\_IN}).
  \item \textbf{Autorisation sur ressources}~: vérification de propriété avant toute
        suppression (photo, commentaire) ou action réservée au créateur du groupe.
  \item \textbf{Logger de requêtes}~: toutes les requêtes entrantes sont journalisées
        (méthode, URL, IP, durée, body avec masquage du champ \texttt{password}).
  \item \textbf{Secrets}~: externalisés dans \texttt{.env} (non versionné) -- URI MongoDB,
        JWT secret, clés Cloudinary, compte de service Firebase.
\end{itemize}

% ─────────────────────────────────────────────────────────────────────────────
\chapter{Fonctionnalités réalisées -- TravelShare}
% ─────────────────────────────────────────────────────────────────────────────

\section{Authentification}

\subsection{Connexion et inscription}

Le module d'authentification est implémenté dans \texttt{AuthDialog}, une boîte de dialogue
modale à deux onglets (Connexion / Inscription).

\begin{itemize}
  \item \textbf{Connexion}~: envoi de l'e-mail et du mot de passe vers \texttt{POST /api/auth/login}.
        En cas de succès, le token JWT et les informations utilisateur (nom complet,
        pseudo, identifiant MongoDB) sont persistés localement via \texttt{TokenManager}
        (SharedPreferences chiffré).
  \item \textbf{Inscription}~: collecte du nom complet, du pseudo, de l'e-mail et du mot de
        passe, validation côté client (champs non vides, longueur minimale du mot de passe),
        puis appel à \texttt{POST /api/auth/register}.
  \item L'interface désactive les boutons pendant l'appel réseau pour prévenir les
        soumissions multiples.
\end{itemize}

\begin{figure}[H]
  \centering
  \begin{tikzpicture}
    \node[draw=medgray, dashed, fill=lightgray, minimum width=5cm, minimum height=3cm,
          font=\small\itshape\color{darkgray}, align=center, rounded corners=4pt]
          {Capture d'écran~:\\[0.3cm]\textbf{Dialogue d'authentification}\\(onglets Connexion / Inscription)};
  \end{tikzpicture}
  \caption{Dialogue d'authentification -- onglets Connexion et Inscription}
  \label{fig:auth}
\end{figure}

\subsection{Gestion de session}

La classe \texttt{TokenManager} centralise la persistance de session~:
\begin{itemize}
  \item Sauvegarde du token JWT, du nom complet, du pseudo, de l'identifiant et de la date
        de création du compte.
  \item Vérification de l'état de connexion (\texttt{isLoggedIn()}) utilisée dans
        l'ensemble de l'application pour adapter dynamiquement l'interface.
  \item Invalidation complète à la déconnexion (\texttt{clear()}).
\end{itemize}

L'état de connexion est exposé via un \texttt{LiveData<Boolean>} dans \texttt{MainViewModel},
ce qui permet à toutes les vues de s'adapter automatiquement (affichage du bouton de
connexion, visibilité des FAB, etc.).

\section{Écran d'accueil -- Fil de photos}

\texttt{HomeFragment} est l'écran central de l'application. Il intègre les éléments suivants~:

\subsection{Chargement et affichage des photos}

Les photos publiques sont chargées depuis \texttt{GET /api/photos} (page~1, limite~100) par
\texttt{PhotoViewModel}. La liste est stockée dans un \texttt{MutableLiveData<List<Photo>>}
et filtrée localement à chaque changement de critère, sans nouvel appel réseau.

Trois modes d'affichage sont disponibles, commutables à la volée~:

\begin{itemize}
  \item \textbf{Vue grille}~: \texttt{GridLayoutManager} 2 colonnes, bannière anonyme
        occupant toute la largeur (span~2).
  \item \textbf{Vue liste}~: \texttt{LinearLayoutManager}, cards avec plus d'informations.
  \item \textbf{Vue carte}~: \texttt{GoogleMap} avec marqueurs positionnés sur les
        coordonnées GPS de chaque photo~; un clic sur le marqueur ouvre un
        \texttt{BottomSheet} de prévisualisation.
\end{itemize}

Un \texttt{SwipeRefreshLayout} permet le rechargement manuel du fil.

\begin{figure}[H]
  \centering
  \begin{subcaptionbox}{Vue grille\label{fig:home_grid}}[0.32\textwidth]
    \begin{tikzpicture}
      \node[draw=medgray, dashed, fill=lightgray, minimum width=3.8cm, minimum height=3cm,
            font=\tiny\itshape\color{darkgray}, align=center, rounded corners=4pt]
            {Capture d'écran\\vue grille};
    \end{tikzpicture}
  \end{subcaptionbox}
  \hfill
  \begin{subcaptionbox}{Vue liste\label{fig:home_list}}[0.32\textwidth]
    \begin{tikzpicture}
      \node[draw=medgray, dashed, fill=lightgray, minimum width=3.8cm, minimum height=3cm,
            font=\tiny\itshape\color{darkgray}, align=center, rounded corners=4pt]
            {Capture d'écran\\vue liste};
    \end{tikzpicture}
  \end{subcaptionbox}
  \hfill
  \begin{subcaptionbox}{Vue carte\label{fig:home_map}}[0.32\textwidth]
    \begin{tikzpicture}
      \node[draw=medgray, dashed, fill=lightgray, minimum width=3.8cm, minimum height=3cm,
            font=\tiny\itshape\color{darkgray}, align=center, rounded corners=4pt]
            {Capture d'écran\\vue carte};
    \end{tikzpicture}
  \end{subcaptionbox}
  \caption{Les trois modes de visualisation des photos}
  \label{fig:home_views}
\end{figure}

\subsection{Recherche textuelle et vocale}

\begin{itemize}
  \item \textbf{Recherche par saisie}~: un champ \texttt{EditText} déclenche un filtre
        local en temps réel (via \texttt{TextWatcher}) sur le nom du lieu, le pays, la
        description, le nom de l'auteur et les tags.
  \item \textbf{Recherche vocale}~: le bouton microphone (icône \texttt{ic\_mic}) lance
        l'\texttt{Intent} \texttt{RecognizerIntent.ACTION\_RECOGNIZE\_SPEECH}. Le texte
        reconnu est injecté directement dans le champ de recherche, déclenchant ainsi le
        filtre. La permission \texttt{RECORD\_AUDIO} est demandée au runtime.
\end{itemize}

\subsection{Filtres avancés}

Le bouton \emph{Filtres} ouvre \texttt{FiltersDialog}, une boîte de dialogue proposant~:

\begin{table}[H]
\centering
\caption{Critères de filtrage disponibles}
\begin{tabularx}{\textwidth}{lX}
\toprule
\rowcolor{primary!15}
\textbf{Critère} & \textbf{Implémentation} \\
\midrule
Type de lieu    & Spinner (Nature, Plage, Urbain, Montagne, Musée, \ldots) \\
Auteur          & Spinner alimenté dynamiquement avec les auteurs des photos chargées \\
Période         & Deux DatePickerDialog (date début / date fin), comparaison avec \texttt{createdAt} \\
Rayon géo.      & Latitude, longitude (saisie ou sélection sur carte via \texttt{MapPickerActivity}) + rayon en km~; calcul de distance Haversine \\
Photo similaire & Saisie d'un identifiant de photo (champ prévu, traitement partiel) \\
\bottomrule
\end{tabularx}
\end{table}

\begin{figure}[H]
  \centering
  \begin{tikzpicture}
    \node[draw=medgray, dashed, fill=lightgray, minimum width=7cm, minimum height=3cm,
          font=\small\itshape\color{darkgray}, align=center, rounded corners=4pt]
          {Capture d'écran~:\\[0.3cm]\textbf{Dialogue des filtres avancés}};
  \end{tikzpicture}
  \caption{Dialogue des filtres avancés}
  \label{fig:filters}
\end{figure}

\section{Fiche détaillée d'une photo}

\texttt{PhotoDetailActivity} affiche l'ensemble des informations associées à une photo~:

\begin{itemize}
  \item \textbf{Image} en haute résolution (chargée via Glide), cliquable pour accéder à
        la vue plein écran avec zoom (\texttt{FullScreenImageActivity} + PhotoView).
  \item \textbf{Métadonnées}~: auteur (avec avatar-initiale), date, type de lieu, localisation
        (ville + pays).
  \item \textbf{Description} et \textbf{how to get there} (comment s'y rendre).
  \item \textbf{Compteur de likes} et bouton Like/Unlike (bascule en rouge/gris).
  \item \textbf{Bouton Partager} : \texttt{Intent.ACTION\_SEND} vers les applications
        tierces installées.
  \item \textbf{Bouton Signaler} : dialogue de confirmation (signalement logique, pas encore
        d'endpoint API dédié).
  \item \textbf{Bouton Itinéraire} (\emph{Ouvrir dans Maps}) : construction d'un URI
        \texttt{geo:} avec coordonnées et nom du lieu, lancement de Google Maps via Intent ;
        repli sur le navigateur web si Maps n'est pas installé.
  \item \textbf{Section commentaires} : chargement asynchrone, ajout et suppression par
        l'auteur (mode connecté uniquement).
\end{itemize}

\begin{figure}[H]
  \centering
  \begin{subcaptionbox}{Fiche photo\label{fig:detail_main}}[0.48\textwidth]
    \begin{tikzpicture}
      \node[draw=medgray, dashed, fill=lightgray, minimum width=5.5cm, minimum height=3cm,
            font=\small\itshape\color{darkgray}, align=center, rounded corners=4pt]
            {Capture d'écran~:\\fiche détaillée de photo};
    \end{tikzpicture}
  \end{subcaptionbox}
  \hfill
  \begin{subcaptionbox}{Plein écran + zoom\label{fig:detail_fullscreen}}[0.48\textwidth]
    \begin{tikzpicture}
      \node[draw=medgray, dashed, fill=lightgray, minimum width=5.5cm, minimum height=3cm,
            font=\small\itshape\color{darkgray}, align=center, rounded corners=4pt]
            {Capture d'écran~:\\affichage plein écran avec zoom};
    \end{tikzpicture}
  \end{subcaptionbox}
  \caption{Fiche détaillée d'une photo et vue plein écran}
  \label{fig:photo_detail}
\end{figure}

\section{Publication de photos}

\texttt{PublishBottomSheet} est accessible via le bouton flottant (FAB) visible uniquement
en mode connecté. Le flux de publication se déroule en deux étapes~:

\begin{enumerate}
  \item \textbf{Upload de l'image}~: sélection depuis la galerie Android, copie dans le
        cache de l'application, envoi multipart vers \texttt{POST /api/photos/upload}
        (stockage Cloudinary). L'URL Cloudinary renvoyée est sauvegardée.
  \item \textbf{Création de la fiche photo}~: appel \texttt{POST /api/photos} avec tous
        les métadonnées saisies par l'utilisateur.
\end{enumerate}

\textbf{Champs du formulaire de publication~:}
\begin{itemize}
  \item Description libre.
  \item Lieu (ville) et pays.
  \item Type de lieu (spinner~: Nature, Urbain, Plage, Montagne, Musée, Magasin, Restaurant).
  \item Tags (ajout par chips avec possibilité de suppression individuelle).
  \item Coordonnées GPS (saisie manuelle ou sélection sur carte via \texttt{MapPickerActivity}).
  \item \emph{Comment s'y rendre} (texte libre).
  \item \textbf{Visibilité}~: Public (accessible à tous) ou Privé (réservé à des groupes
        sélectionnés via des cases à cocher).
\end{itemize}

\begin{figure}[H]
  \centering
  \begin{subcaptionbox}{Formulaire de publication\label{fig:publish_form}}[0.48\textwidth]
    \begin{tikzpicture}
      \node[draw=medgray, dashed, fill=lightgray, minimum width=5.5cm, minimum height=3cm,
            font=\small\itshape\color{darkgray}, align=center, rounded corners=4pt]
            {Capture d'écran~:\\bottom sheet de publication};
    \end{tikzpicture}
  \end{subcaptionbox}
  \hfill
  \begin{subcaptionbox}{Sélection de lieu sur carte\label{fig:map_picker}}[0.48\textwidth]
    \begin{tikzpicture}
      \node[draw=medgray, dashed, fill=lightgray, minimum width=5.5cm, minimum height=3cm,
            font=\small\itshape\color{darkgray}, align=center, rounded corners=4pt]
            {Capture d'écran~:\\MapPickerActivity};
    \end{tikzpicture}
  \end{subcaptionbox}
  \caption{Publication d'une photo et sélection du lieu}
  \label{fig:publish}
\end{figure}

\section{Gestion des groupes}

\texttt{GroupsBottomSheet} propose deux onglets~:

\paragraph{Mes groupes}
Affichage des groupes dont l'utilisateur est membre. Un appui long ouvre un menu contextuel
permettant de \emph{quitter} le groupe ou, si l'utilisateur en est le créateur, de le
\emph{supprimer}. La création d'un nouveau groupe (nom, description, image de couverture)
est accessible depuis cet onglet.

\paragraph{Découvrir}
Liste des groupes existants auxquels l'utilisateur n'appartient pas encore, avec un bouton
\emph{Rejoindre}.

Un clic sur un groupe ouvre \texttt{GroupPhotosActivity}, qui affiche les photos partagées
au sein de ce groupe.

\begin{figure}[H]
  \centering
  \begin{subcaptionbox}{Mes groupes\label{fig:groups_my}}[0.32\textwidth]
    \begin{tikzpicture}
      \node[draw=medgray, dashed, fill=lightgray, minimum width=3.8cm, minimum height=3cm,
            font=\tiny\itshape\color{darkgray}, align=center, rounded corners=4pt]
            {Mes groupes};
    \end{tikzpicture}
  \end{subcaptionbox}
  \hfill
  \begin{subcaptionbox}{Découvrir des groupes\label{fig:groups_discover}}[0.32\textwidth]
    \begin{tikzpicture}
      \node[draw=medgray, dashed, fill=lightgray, minimum width=3.8cm, minimum height=3cm,
            font=\tiny\itshape\color{darkgray}, align=center, rounded corners=4pt]
            {Découvrir des groupes};
    \end{tikzpicture}
  \end{subcaptionbox}
  \hfill
  \begin{subcaptionbox}{Photos d'un groupe\label{fig:group_photos}}[0.32\textwidth]
    \begin{tikzpicture}
      \node[draw=medgray, dashed, fill=lightgray, minimum width=3.8cm, minimum height=3cm,
            font=\tiny\itshape\color{darkgray}, align=center, rounded corners=4pt]
            {Photos du groupe};
    \end{tikzpicture}
  \end{subcaptionbox}
  \caption{Écrans de gestion des groupes}
  \label{fig:groups}
\end{figure}

\section{Notifications}

\subsection{Notifications push (Firebase Cloud Messaging)}

Le service \texttt{MyFirebaseMessagingService} étend \texttt{FirebaseMessagingService}~:
\begin{itemize}
  \item À chaque renouvellement du token FCM, ce dernier est envoyé au backend via
        \texttt{PUT /api/auth/fcm-token} (si l'utilisateur est connecté).
  \item À la réception d'un message push, une notification système est créée avec
        \texttt{NotificationCompat.Builder} (canal haute priorité, redirection vers
        \texttt{MainActivity}).
\end{itemize}

\subsection{Centre de notifications in-app}

\texttt{NotificationsBottomSheet} affiche l'historique des notifications en base~:
\begin{itemize}
  \item Chargement depuis \texttt{GET /api/notifications} et \texttt{GET /api/notifications/unread-count}.
  \item Un badge numérique sur l'icône de la barre de navigation indique le nombre
        de notifications non lues (observable via \texttt{MainViewModel}).
  \item Un clic sur une notification la marque comme lue (\texttt{PATCH /api/notifications/\{id\}/read})
        et redirige vers la ressource concernée (photo ou groupe).
\end{itemize}

\begin{figure}[H]
  \centering
  \begin{tikzpicture}
    \node[draw=medgray, dashed, fill=lightgray, minimum width=6cm, minimum height=3cm,
          font=\small\itshape\color{darkgray}, align=center, rounded corners=4pt]
          {Capture d'écran~:\\[0.3cm]\textbf{Centre de notifications}\\(badge + liste)};
  \end{tikzpicture}
  \caption{Centre de notifications in-app avec badge}
  \label{fig:notifications}
\end{figure}

\section{Profil utilisateur}

\texttt{ProfileFragment} s'adapte à l'état de connexion~:
\begin{itemize}
  \item \textbf{Non connecté}~: message d'invitation avec bouton de connexion.
  \item \textbf{Connecté}~: appel à \texttt{GET /api/auth/me} pour obtenir les données
        fraîches depuis le serveur (nom complet, pseudo, date de création du compte).
        Affichage de l'avatar-initiale et de l'année d'inscription.
        Bouton de déconnexion.
\end{itemize}

\begin{figure}[H]
  \centering
  \begin{subcaptionbox}{Non connecté\label{fig:profile_out}}[0.48\textwidth]
    \begin{tikzpicture}
      \node[draw=medgray, dashed, fill=lightgray, minimum width=5.5cm, minimum height=3cm,
            font=\small\itshape\color{darkgray}, align=center, rounded corners=4pt]
            {Capture d'écran~:\\profil non connecté};
    \end{tikzpicture}
  \end{subcaptionbox}
  \hfill
  \begin{subcaptionbox}{Connecté\label{fig:profile_in}}[0.48\textwidth]
    \begin{tikzpicture}
      \node[draw=medgray, dashed, fill=lightgray, minimum width=5.5cm, minimum height=3cm,
            font=\small\itshape\color{darkgray}, align=center, rounded corners=4pt]
            {Capture d'écran~:\\profil utilisateur connecté};
    \end{tikzpicture}
  \end{subcaptionbox}
  \caption{Fragment de profil utilisateur}
  \label{fig:profile}
\end{figure}

\section{Modèle de données}

Le modèle \texttt{Photo} est le cœur applicatif. Il expose~:

\begin{lstlisting}[language=Java, caption=Extrait du modèle Photo.java]
public class Photo {
    @SerializedName("_id")   private String mongoId;
    private String imageUrl, location, country, locationType;
    private double latitude, longitude;
    private String description, howToGetThere;
    @SerializedName("createdAt") private String date;
    @SerializedName("likes")     private List<String> likesArray;
    private boolean likedByMe, isPublic;
    private List<String> tags;
    private Author author;   // sous-objet { _id, fullName, username, avatar }
}
\end{lstlisting}

% ─────────────────────────────────────────────────────────────────────────────
\chapter{Tableau de bord des fonctionnalités}
% ─────────────────────────────────────────────────────────────────────────────

Le tableau ci-dessous synthétise l'état d'avancement de chaque fonctionnalité spécifiée.

\begin{table}[H]
\centering
\caption{État d'avancement des fonctionnalités TravelShare}
\renewcommand{\arraystretch}{1.4}
\begin{tabularx}{\textwidth}{lXc}
\toprule
\rowcolor{primary!20}
\textbf{Catégorie} & \textbf{Fonctionnalité} & \textbf{État} \\
\midrule
\multirow{4}{*}{\textbf{Authentification}}
  & Inscription (nom, pseudo, e-mail, mot de passe) & \cmark \\
  & Connexion par e-mail / mot de passe              & \cmark \\
  & Persistance de session (JWT + SharedPreferences) & \cmark \\
  & Déconnexion                                      & \cmark \\
\midrule
\multirow{5}{*}{\textbf{Parcours photos}}
  & Chargement du fil public (API REST)              & \cmark \\
  & Vue grille                                       & \cmark \\
  & Vue liste                                        & \cmark \\
  & Vue carte (Google Maps + marqueurs)              & \cmark \\
  & Pull-to-refresh                                  & \cmark \\
\midrule
\multirow{6}{*}{\textbf{Recherche \& filtres}}
  & Recherche textuelle (temps réel)                 & \cmark \\
  & Recherche vocale (RecognizerIntent)              & \cmark \\
  & Filtre par type de lieu                          & \cmark \\
  & Filtre par auteur                                & \cmark \\
  & Filtre par période (date début / fin)            & \cmark \\
  & Filtre par rayon géographique (Haversine)        & \cmark \\
  & Filtre par similarité (photos semblables)        & \pmark \\
\midrule
\multirow{6}{*}{\textbf{Fiche photo}}
  & Affichage image + métadonnées complètes          & \cmark \\
  & Vue plein écran avec zoom (PhotoView)            & \cmark \\
  & Like / Unlike                                    & \cmark \\
  & Commentaires (lecture, ajout, suppression)       & \cmark \\
  & Partage via Intent Android                       & \cmark \\
  & Signalement (dialogue UI)                        & \pmark \\
  & Itinéraire vers le lieu (Intent Google Maps)     & \cmark \\
\midrule
\multirow{5}{*}{\textbf{Publication}}
  & Sélection photo depuis la galerie                & \cmark \\
  & Upload vers Cloudinary (multipart)               & \cmark \\
  & Formulaire complet (lieu, tags, coordonnées…)    & \cmark \\
  & Sélection du lieu sur carte (MapPickerActivity)  & \cmark \\
  & Visibilité public / groupes                      & \cmark \\
  & Suppression de sa propre photo                   & \cmark \\
  & Annotations assistées par IA                     & \xmark \\
\midrule
\multirow{5}{*}{\textbf{Groupes}}
  & Création d'un groupe                             & \cmark \\
  & Rejoindre / quitter un groupe                    & \cmark \\
  & Supprimer un groupe (créateur)                   & \cmark \\
  & Découverte des groupes publics                   & \cmark \\
  & Photos partagées dans un groupe                  & \cmark \\
\midrule
\multirow{4}{*}{\textbf{Notifications}}
  & Notifications push FCM                           & \cmark \\
  & Centre de notifications in-app                   & \cmark \\
  & Badge sur l'onglet de navigation                 & \cmark \\
  & Configuration des abonnements (thème/auteur)     & \xmark \\
\midrule
\multirow{3}{*}{\textbf{Profil}}
  & Affichage des informations utilisateur           & \cmark \\
  & Galerie des photos publiées par l'utilisateur    & \xmark \\
  & Modification du profil                           & \xmark \\
\bottomrule
\multicolumn{3}{l}{\footnotesize \cmark~= Réalisé \quad \pmark~= Partiellement réalisé \quad \xmark~= Non encore réalisé}
\end{tabularx}
\end{table}

% ─────────────────────────────────────────────────────────────────────────────
\chapter{Points techniques remarquables}
% ─────────────────────────────────────────────────────────────────────────────

\section{Moteur de filtrage local (client-side)}

Pour éviter des allers-retours réseau à chaque frappe, tous les filtres sont appliqués
localement sur la liste complète des photos déjà chargée. La classe \texttt{FilterCriteria}
encapsule l'ensemble des critères actifs. La méthode \texttt{applyFilters()} du
\texttt{PhotoViewModel} itère sur cette liste et applique en cascade les critères de type,
d'auteur, de période (comparaison \texttt{Date}), de rayon (formule Haversine) et de texte
(correspondance dans lieu, pays, description, auteur, tags).

\section{Calcul de distance Haversine}

La formule de Haversine est utilisée pour calculer la distance en kilomètres entre deux
points GPS, ce qui permet le filtre \emph{autour d'un lieu} sans dépendance à un service
tiers~:

\begin{lstlisting}[language=Java, caption=Calcul de distance Haversine (PhotoViewModel)]
private double calculateDistance(double lat1, double lng1,
                                 double lat2, double lng2) {
    double R = 6371; // rayon moyen de la Terre en km
    double dLat = Math.toRadians(lat2 - lat1);
    double dLng = Math.toRadians(lng2 - lng1);
    double a = Math.sin(dLat/2)*Math.sin(dLat/2)
             + Math.cos(Math.toRadians(lat1))
             * Math.cos(Math.toRadians(lat2))
             * Math.sin(dLng/2)*Math.sin(dLng/2);
    double c = 2 * Math.atan2(Math.sqrt(a), Math.sqrt(1-a));
    return R * c;
}
\end{lstlisting}

\section{Configuration dynamique de l'URL du serveur}

L'écran \texttt{IpConfigActivity} permet de modifier à la volée l'adresse IP du backend
REST. Cette fonctionnalité est particulièrement utile lors des phases de test en réseau
local (e.g.\ machine virtuelle ou appareil physique sur le même sous-réseau Wi-Fi). La
nouvelle URL est passée à \texttt{RetrofitClient.setBaseUrl()}, ce qui force la
réinstanciation du client Retrofit.

\section{Adaptateur multi-vues avec bannière anonyme}

\texttt{PhotoAdapter} gère deux types d'items (\texttt{TYPE\_BANNER} et \texttt{TYPE\_PHOTO})
et trois modes d'affichage (grille, liste, card). La bannière est injectée en première
position pour inviter les utilisateurs non connectés à créer un compte. La
\texttt{GridLayoutManager.SpanSizeLookup} lui alloue automatiquement une largeur pleine
(span~= 2).

% ─────────────────────────────────────────────────────────────────────────────
\chapter{Fonctionnalités restantes et plan de travail}
% ─────────────────────────────────────────────────────────────────────────────

\section{Fonctionnalités prioritaires à compléter}

\begin{warningbox}[Travaux restants -- priorité haute]
\begin{enumerate}[leftmargin=*]
  \item \textbf{Annotations assistées par IA}~: enrichissement automatique de la photo
        lors de la publication (génération de tags, description courte) via un appel à
        une API d'IA (Claude, Gemini ou Vision API). C'est une fonctionnalité explicitement
        mentionnée dans la spécification.
  \item \textbf{Galerie de profil}~: affichage des photos publiées par l'utilisateur
        connecté dans son onglet Profil, avec possibilité de gestion (suppression).
  \item \textbf{Configuration des abonnements de notification}~: permettre à l'utilisateur
        de choisir les événements qui le notifient (publication d'un auteur donné, d'un
        groupe, d'un lieu ou d'un thème).
  \item \textbf{Signalement côté API}~: implémenter l'endpoint de signalement et connecter
        le bouton Signaler de \texttt{PhotoDetailActivity} à cet endpoint.
  \item \textbf{Filtre par similarité}~: la logique de récupération des photos similaires
        à une photo de référence (champ \texttt{similarPhotoId} déjà présent dans
        \texttt{FilterCriteria}) doit être branchée côté ViewModel et API.
\end{enumerate}
\end{warningbox}

\section{Améliorations envisagées}

\begin{itemize}
  \item \textbf{Pagination incrémentale}~: remplacer le chargement en une seule page
        (limite fixe à 100) par un défilement infini (\texttt{RecyclerView} +
        \texttt{EndlessScrollListener}).
  \item \textbf{Modification du profil}~: formulaire d'édition du nom, du pseudo et de
        l'avatar.
  \item \textbf{Mode hors-ligne partiel}~: mise en cache des dernières photos consultées
        (Room ou fichiers JSON dans le stockage interne).
  \item \textbf{Intégration TravelPath (Partie~3)}~: ajout du module passerelle permettant
        de naviguer entre TravelShare et TravelPath dans une navigation unifiée.
\end{itemize}

\section{Estimation de l'avancement global}

\begin{figure}[H]
\centering
\begin{tikzpicture}
  \def\ptotal{15}
  \def\pdone{11}
  \def\ppartial{2}
  \def\barW{12}
  \fill[lightgray, rounded corners=4pt] (0,0) rectangle (\barW, 0.8);
  \fill[success!70, rounded corners=4pt] (0,0) rectangle (\barW*\pdone/\ptotal, 0.8);
  \fill[warning!80] (\barW*\pdone/\ptotal, 0) rectangle (\barW*(\pdone+\ppartial)/\ptotal, 0.8);
  \draw[medgray, rounded corners=4pt] (0,0) rectangle (\barW, 0.8);
  \node[anchor=west, font=\small\bfseries] at (0.2, 0.4) {\textcolor{white}{Réalisé}};
  \node[font=\small\bfseries] at (\barW*(\pdone+0.5*\ppartial)/\ptotal, 0.4) {\textcolor{white}{\pmark}};
  \node[right, font=\small, color=darkgray] at (\barW + 0.2, 0.4)
        {\textbf{\pdone/\ptotal} fonctionnalités};
\end{tikzpicture}
\caption{Estimation visuelle de l'avancement (fonctionnalités principales)}
\label{fig:progress}
\end{figure}

Sur les \textbf{15~fonctionnalités principales} identifiées dans la spécification~:
\begin{itemize}
  \item \textbf{11} sont entièrement réalisées (\cmark),
  \item \textbf{2} sont partiellement réalisées (\pmark) : filtre par similarité et signalement,
  \item \textbf{2} restent à réaliser (\xmark) : annotations IA et configuration des abonnements de notification.
\end{itemize}

% ─────────────────────────────────────────────────────────────────────────────
\chapter{Partie 2 : TravelPath}
% ─────────────────────────────────────────────────────────────────────────────

\noindent\textbf{Dépôt GitLab :}
\url{https://gitlab.etu.umontpellier.fr/e20190009868/travelsharepath}

\section*{Stack technique}

\begin{itemize}
  \item \textbf{Frontend} : Android natif en Java
  \item \textbf{Backend} : Node.js / Express / MongoDB
  \item \textbf{Architecture} : MVVM (ViewModel + LiveData), Retrofit2, Navigation Component
\end{itemize}

\section{Authentification}

\begin{itemize}
  \item Inscription avec nom complet, email et mot de passe .
  \item Connexion avec un token JWT persisté dans les SharedPreferences.
  \item Déconnexion avec nettoyage complet de la session.
  \item L'écran d'accueil adapte son interface selon l'état de connexion.
\end{itemize}

\section{Création d'un parcours}

Le formulaire de création est organisé en deux niveaux.

\textbf{Niveau parcours :}
\begin{itemize}
  \item Titre, description, ville.
  \item Météo : Ensoleillé, Nuageux, Pluvieux, Tout temps.
  \item Niveau d'effort : Facile, Moyen (Actif), Difficile.
\end{itemize}

\textbf{Niveau étape (ajout multiple) :}
\begin{itemize}
  \item Nom du lieu, géocodé automatiquement via Nominatim/OSM.
  \item Type : Culture, Découverte, Restauration, Loisirs.
  \item Durée (minutes), budget (€), description.
  \item Photos multiples ajoutées depuis la galerie.
  \item Réordonnancement par glisser-déposer.
  \item Édition inline de chaque étape.
  \item Compteurs en temps réel (durée totale, budget total, nombre d'étapes).
\end{itemize}

Un nouveau parcours est créé en statut \textit{Brouillon}.

\section{Recherche de parcours}

\begin{itemize}
  \item Formulaire de filtres : ville, budget (slider), durée (slider), niveau d'effort,
    types d'activités (multi-sélection).
  \item Seuls les parcours \textit{publiés} apparaissent dans les résultats.
  \item Filtrage côté serveur (ville, budget max, durée max, effort, types d'activités, météo).
\end{itemize}

\section{Interactions sociales}

\textbf{Aimer un parcours} :
\begin{itemize}
  \item Bouton cœur avec compteur visible sur la fiche détail et dans les listes.
  \item Toggle like/unlike confirmé par le serveur (pas d'optimistic update).
  \item Le compteur se met à jour après confirmation serveur.
\end{itemize}

\textbf{Mettre en favori} :
\begin{itemize}
  \item Icône étoile sur la fiche détail.
  \item Liste « Mes Favoris » dédiée avec les mêmes actions disponibles
    (like, dé-favori, régénérer).
  \item État favori synchronisé via \texttt{GET /auth/me}.
\end{itemize}

\section{Publication d'un parcours}

\begin{itemize}
  \item Bouton « Publier » visible uniquement sur les parcours en statut Brouillon.
  \item Passage à statut \textit{Publié} côté serveur.
  \item La liste « Mes parcours » se rafraîchit automatiquement après publication.
\end{itemize}

\section{Détail d'un parcours --- vue Itinéraire}

\textbf{En-tête :}
\begin{itemize}
  \item Titre, budget total, durée totale, distance estimée, effort, météo, date.
  \item Jusqu'à 4 photos d'aperçu (1\textsuperscript{re} photo de chaque étape).
\end{itemize}

\textbf{Par étape :}
\begin{itemize}
  \item Numéro et nom du lieu.
  \item Badge de type coloré : Culture (indigo), Restauration (orange),
    Loisirs (vert), Découverte (bleu).
  \item Description, durée, budget.
  \item Niveau d'effort coloré (vert / orange / rouge).
  \item Galerie photos swipeable horizontalement.
  \item Temps de marche vers l'étape suivante (si renseigné).
\end{itemize}

\section{Détail d'un parcours --- vue Carte}

\begin{itemize}
  \item Carte Leaflet.js intégrée dans une WebView, avec tuiles CartoDB.
  \item Marqueurs numérotés (cercles rouges) pour chaque étape.
  \item Géocodage à la volée des étapes sans coordonnées GPS.
  \item Tracé d'abord en pointillés bleus, puis remplacé par l'itinéraire routier
    réel calculé via OSRM.
\end{itemize}

\section{Autres fonctionnalités}

\begin{tabularx}{\linewidth}{@{}>{\bfseries}lX@{}}
\toprule
Mes Parcours & Liste avec actions : voir, éditer, publier, supprimer (avec confirmation). \\
\midrule
Édition / Régénération & Propriétaire $\rightarrow$ mise à jour ; autre utilisateur
  $\rightarrow$ fork (copie). \\
\midrule
Export PDF & PDF multi-pages stylé : couverture, statistiques, description, une carte
  par étape. \\
\midrule
Upload photos & Multipart vers le backend, stocké dans \texttt{/backend/uploads/}. \\
\midrule
Profil & Nom, email, statistiques (parcours créés / publiés), lien vers création. \\
\bottomrule
\end{tabularx}

% ─────────────────────────────────────────────────────────────────────────────
\chapter{Passerelle Traveling}
% ─────────────────────────────────────────────────────────────────────────────

\textit{À discuter et à compléter en binôme.}

L'objectif de la passerelle est de faire communiquer les deux applications afin
qu'elles forment un système cohérent. L'architecture envisagée s'appuie sur les
mécanismes de communication inter-applications d'Android :

\begin{itemize}
  \item \textbf{Intents} : navigation entre écrans des deux applications.
  \item \textbf{Content Providers} : exposition de données structurées
    (photos d'un lieu, étapes d'un parcours).
  \item \textbf{Services AIDL} : appels de méthodes distantes pour les calculs
    (par exemple, génération d'un parcours à partir des photos populaires d'un lieu).
\end{itemize}

\subsection*{Scénarios envisagés}

\begin{itemize}
  \item Depuis une photo dans TravelShare $\rightarrow$ générer un parcours TravelPath
    incluant ce lieu.
  \item Depuis une étape d'un parcours TravelPath $\rightarrow$ consulter les photos
    associées au lieu sur TravelShare.
\end{itemize}

% ─────────────────────────────────────────────────────────────────────────────
\chapter{Conclusion}
% ─────────────────────────────────────────────────────────────────────────────

À mi-parcours, la partie \textbf{TravelShare} du projet \emph{Traveling} présente une base
fonctionnelle solide, aussi bien côté \textbf{frontend Android} que côté \textbf{backend REST}.
L'architecture MVVM du client Android garantit une séparation claire des responsabilités et
une bonne maintenabilité. Le serveur Node.js/Express expose l'intégralité des endpoints
nécessaires, persistés en MongoDB Atlas, avec gestion des images (Cloudinary) et des
notifications push (Firebase Cloud Messaging). Les fonctionnalités cœur de l'application
-- parcours du fil de photos, recherche multicritère, publication, groupes et notifications
push -- sont toutes opérationnelles de bout en bout.

Les prochaines semaines seront consacrées principalement à l'intégration des annotations
assistées par IA, à l'enrichissement du profil utilisateur, à la finalisation des
fonctionnalités partielles (signalement, similarité) et à la préparation de la passerelle
d'intégration avec TravelPath.

\begin{successbox}[Points forts à mi-parcours]
\begin{itemize}[noitemsep]
  \item Architecture MVVM propre et extensible (frontend Android).
  \item API REST complète~: 24 endpoints couvrant auth, photos, groupes et notifications.
  \item Index géospatial MongoDB (\texttt{2dsphere}) pour les requêtes de proximité.
  \item Recherche vocale native (RecognizerIntent) fonctionnelle.
  \item Moteur de filtrage multi-critères entièrement côté client (réactif, sans latence réseau).
  \item Upload d'images en multipart streamé directement vers Cloudinary.
  \item Notifications push Firebase Cloud Messaging via API HTTP v1 (sans SDK tiers).
  \item Vue carte Google Maps avec marqueurs cliquables.
  \item Itinéraire depuis la fiche photo vers Google Maps.
\end{itemize}
\end{successbox}

\end{document}
